\scanpage{105}
\begin{center}- \mspage{170} -\end{center}
\margin{IX}
germe d'homomorphisme de fibrés vectoriels $E\longrightarrow E'$, appliquant $R$ dans $R'$, revient à la donnée d'un système cohérent d'homomorphismes $E_n\longrightarrow E'_n$ appliquant $R_n$ dans $R'_n$. Or si $\underline{\mathrm O}_s$ désigne l'anneau local de $S$ en $s$, la donnée d'un système cohérent d'homomorphismes $E_n\longrightarrow E'_n$ revient à la donnée d'un homomorphisme de $\widehat{\underline{\mathrm O}}_s$-modules
\[
\widehat F:\widehat{\underline E}_s\longrightarrow\widehat{\underline E'}_s
\]
(où $\underline E_s$ et $E'_s$ sont les $\underline{\mathrm O}_s$-modules libres fibres en $s$ des $\underline{\mathrm O}_S$-Modules $\underline E$, $\underline E'$ correspondant à $E,E'$), tandis qu'un germe d'homomorphisme de fibrés vectoriels $E\longrightarrow E'$ revient à la donnée d'un homomorphisme de $\underline{\mathrm O}_s$-modules
\[
F:\underline E_s\longrightarrow\underline E'_s\quad.
\]
Cela montre déjà que l'application (13.2.1.1) est injective, et pour la surjectivité, il suffit de prouver que si on a un homomorphisme $\widehat F$ comme dessus, tel que pour tout $n$ l'homomorphisme induit $F_n:E_n\longrightarrow E'_n$ applique $R_n$ dans $R'_n$, alors $\widehat F$ provient d'un $F$ i.e. applique $\underline E_s$ dans $\underline E'_s$, et le germe correspondant d'homomorphismes $E\longrightarrow E'$ applique $R$ dans $R'$. Or, comme par hypothèse c), $R_s=R_0$ engendre $E_s$, il s'ensuit par Nakayama que le groupe $R_0$ (identifié au groupe des sections de $\widehat R$ sur $\widehat S$, ou,au choix,au groupe des germes de sections de $R$ sur $S$) engendre $\widehat E_s$ ; donc, pour prouver que $\widehat F(\underline E_s)\subset\underline E'_s$,il suffit de prouver que $\widehat F(R_0)\subset\underline E'_s$. Or on a $\widehat F(R_0)\subset R'_0$, où $R'_0$ est lui-même interprété comme le groupe des germes de sections de $R'\subset E'$ au voisinage de $s$, d'où l'inclusion annoncée, et par suite l'existence de $F$. De plus on voit que $F(R_0)\subset R'_0$, où $R_0$ et $R'_0$ sont considérés comme des groupes de germes de sections de $E$ resp. $E'$. En vertu de b), $R_0$

\margin{\diagram{\begin{tikzcd}
\mathrm{DR}(A^{1,0},T)\arrow[r]\arrow[d]&A^{**}\\
\mathrm{DR}(\unsure{\Phi_*}\otimes_{\mathbb Z}A_0,T)&\\
{}\arrow[u]&
\end{tikzcd}}}
\margin{488}
\margin{488}
\margin{104}
